\subsection{Complete-Graph Scheduling and Tensor Lifetimes}
\label{sec:graph-schedule}
The sequencer controls all 71 operations independently of the compute array.
Each entry selects an operator, its source and destination buffers, and the
relevant parameter set. Convolution and transposed-convolution entries select
one of 38 weighted-layer descriptors. ReLU, add, and interpolation entries
instead supply their tensor extents and fixed-point parameters to their
respective engines.

Figure~\ref{fig:graph-execution} illustrates this organization. Logical
feature buffers preserve tensors needed by later residual or fusion
operations; scratch buffers hold transient outputs. Their lifetimes follow
the graph rather than a simple alternation after every operation. A buffer
may be reused only after its last consumer completes. Engine completion must
also cover output-memory acceptance before the next dependent operation is
issued.

\begin{figure}[h]
  \centering
  \begin{tikzpicture}[x=1cm,y=1cm,>=Latex,
  box/.style={draw=paperblue,fill=paperlight,rounded corners=2pt,
    align=center,font=\footnotesize,minimum height=.7cm},
  arr/.style={->,draw=paperblue,thick}]
  \node[box,text width=6.9cm] (seq) at (0,0) {71-entry graph sequencer\\
    opcode, parameters, source / destination, completion};
  \node[box,text width=3.1cm] (weighted) at (-1.9,-1.25) {Weighted engine\\
    Conv / transposed conv};
  \node[box,text width=3.1cm] (aux) at (1.9,-1.25) {Auxiliary engines\\
    ReLU / Add / Bilinear};
  \node[box,text width=6.9cm] (mem) at (0,-2.5) {Shared memory access path};
  \node[box,text width=3.1cm] (feat) at (-1.9,-3.65) {Feature buffers\\
    retain live graph values};
  \node[box,text width=3.1cm] (scratch) at (1.9,-3.65) {Scratch buffers\\
    temporary results};
  \draw[arr] (seq.south)--++(0,-.18)-|(weighted.north);
  \draw[arr] (seq.south)--++(0,-.18)-|(aux.north);
  \draw[arr] (weighted.south)--++(0,-.2)-|(mem.north);
  \draw[arr] (aux.south)--++(0,-.2)-|(mem.north);
  \draw[arr] (mem.south)--++(0,-.18)-|(feat.north);
  \draw[arr] (mem.south)--++(0,-.18)-|(scratch.north);
\end{tikzpicture}

  \caption{Execution organization of the 71-operation decoder. The sequencer
  selects one weighted or auxiliary engine for each graph step. Persistent
  feature buffers retain values for later consumers, while scratch buffers
  hold temporary results. This is a logical execution diagram, not a claim
  that all graph edges remain on chip.}
  \Description{A sequencer dispatching to weighted and auxiliary engines with
  a shared memory interface to persistent and scratch buffers.}
  \label{fig:graph-execution}
\end{figure}

\subsection{Convolution, Upsampling, and Auxiliary Operators}
The weighted engine supports the decoder's $1\times1$ and $3\times3$
convolutions together with kernel-2, stride-2 transposed convolution. Layer
descriptors select the shape and address-generation behavior while reusing
the physical array. The final manuscript will pair the transposed-convolution
address mapping with a small worked tile example from the implementation.

The auxiliary engines implement ReLU, residual addition, and bilinear
interpolation. They share the memory access path with the weighted engine.
Each addition converts its two operands to the
destination scale; interpolation similarly requires a specified coordinate
and rounding convention. These operators are included in decoder timing,
even though they do not use the systolic array.

\subsection{INT8 Datapath and Output Representation}
\label{sec:quantization}
Decoder quantization is an enabling implementation step. The encoder retains
its separately specified execution precision. The hardware consumes signed
INT8 operands, retains INT32 partial sums, and uses integer bias and scaling
parameters for output conversion. This follows the general integer-inference
principle of separating narrow operands from wider accumulation
~\cite{jacob2018quantization}.

For a layer output, the numerical contract can be written as
\begin{equation}
  q_y=\operatorname{sat}_{b}\!\left(
      \operatorname{round}\!\left(M(p+b_q)\right)\right),
  \label{eq:requant}
\end{equation}
where $p$ is the integer dot product, $b_q$ is the integer bias, and $M$
represents the exported fixed-point scaling operation. The exact multiplier,
shift, rounding, and clipping rules must match the reference implementation.
Equation~\ref{eq:requant} expresses the contract and does not replace its
bit-level definition. Intermediate outputs use INT8; the final accelerator
output is signed INT16.

\begin{table}[h]
  \caption{Numerical and memory contracts to be fixed for the final
  evaluation. Pending entries require the corresponding export or test record.}
  \label{tab:contracts}
  \centering\small
  \begin{tabular}{p{0.43\columnwidth}p{0.48\columnwidth}}
    \toprule
    Item & Current contract / required detail\\
    \midrule
    Array operands & Signed INT8\\
    Partial sums & INT32\\
    Intermediate / final output & INT8 / signed INT16\\
    Bias and output scaling & Exported integer parameters\\
    Calibration procedure & \pending\\
    Rounding / saturation tests & \pending\\
    Depth unit and output scale & \pending\\
    \bottomrule
  \end{tabular}
\end{table}

\subsection{Implementation and Verification Boundary}
The accelerator combines RTL control and weighted computation with HLS-based
auxiliary engines. Verification proceeds from the integer software reference
to operator-level RTL checks, complete-graph simulation, and board-level
execution. Matching a software model or generating a bitstream alone does not
establish complete hardware correctness. The final resource table will use
the evaluated bitstream's post-fit report and identify the clock domains for
which timing is met.

\draftnote{Add the author's detailed PE/buffer inset to the main architecture
figure if available. Quantization settings and measured resource counts are
intentionally not inferred from incomplete integration tests.}
